\section{Results}
\label{sec:results}

We present results in three parts corresponding to our research questions:
framework validation (RQ1), the preliminary SEP finding (RQ2), and the
checkpoint aliasing confound (RQ3).

\subsection{Framework Validation (RQ1): End-to-End Infrastructure Correctness}

The measurement framework executes correctly end-to-end. All five modules
(ast\_decomposition, ast\_metric, kl\_metric, sep\_analysis, statistical\_tests)
run without errors on real DeepSeek-Coder-7B-instruct-v1.5 checkpoints.

\textbf{Proof-of-concept results (h-e1, synthetic data):}

\begin{table}[h]
\centering
\caption{Proof-of-concept results on synthetic h-e1 data}
\label{tab:poc_results}
\begin{tabular}{lll}
\toprule
\textbf{Metric} & \textbf{GRPO} & \textbf{DPO} \\
\midrule
Mean semantic AST edit distance & 3.500 & 1.000 \\
Mean edit-per-KL & $\sim$25.9 (low-KL pairs) & $\sim$3.7 (low-KL pairs) \\
Syntax pass rate & 6/6 (1.000) & 6/6 (1.000) \\
Bootstrap 95\% CI (differential) & [4.6500, 8.7314] & --- \\
Mean differential & 6.5047 & --- \\
\bottomrule
\end{tabular}
\end{table}

Figure~\ref{fig:edit_per_kl} shows the semantic-edit-per-KL comparison across the 27 KL-matched
pairs on the proof-of-concept data. The CI excludes zero (lower bound 4.65 $>$ 0),
confirming the measurement infrastructure correctly detects a difference when
one is engineered to exist. Figure~\ref{fig:bootstrap_ci} shows the bootstrap distribution of the
mean differential, with the 95\% CI clearly separated from zero. Figure~\ref{fig:ast_edit_distances}
shows per-problem AST edit distances: GRPO produces higher or equal distances on 4 of 6 problems.

\footnotetext{The above results are from h-e1 synthetic (hand-crafted) proof-of-concept data,
not real GRPO training outputs. DPO training in h-e1 used stub preference pairs rather than
genuine execution-oracle pairs (see Limitation L4 in Section~\ref{sec:limitations}), which may
suppress DPO's structural activity and inflate the apparent raw edit distance advantage.}

\textbf{Component-level verification:}

\begin{table}[h]
\centering
\caption{Component verification status}
\label{tab:components}
\begin{tabular}{lll}
\toprule
\textbf{Component} & \textbf{Status} & \textbf{Evidence} \\
\midrule
FA-AST taxonomy (CF+DF node classification) & \checkmark Functional & SEP values in valid [0, 1] range \\
ZSS semantic edit distance & \checkmark Functional & Correct distances on 6 problems \\
KL log matching (tolerance=0.15) & \checkmark Functional & 27 matched pairs found \\
Bootstrap CI (10,000 samples) & \checkmark Functional & CI correctly excludes zero on PoC \\
Mann-Whitney U test & \checkmark Functional & Test executes; requires unique checkpoints \\
\bottomrule
\end{tabular}
\end{table}

The framework is not merely theoretical --- it is a validated, executable
measurement tool ready for deployment on real training runs with sufficient
checkpoint diversity.

\subsection{Preliminary SEP Finding (RQ2): Near-Equal Semantic Edit Proportions}

\textit{Note: The raw edit distance results in Section~\ref{sec:results}.1 use synthetic h-e1
proof-of-concept data; the SEP results here use real checkpoints from h-m1 but are
underpowered ($n_\text{eff}\approx2$ due to checkpoint aliasing, Section~\ref{sec:aliasing}).
These experiments address different questions and cannot be directly compared.}

The mechanism hypothesis --- that execution reward selectively concentrates
policy movement on semantic AST nodes --- is not supported by the preliminary
analysis. Table~\ref{tab:sep_results} shows the SEP results from h-m1.

\begin{table}[h]
\centering
\caption{Semantic Edit Proportion (SEP) --- Preliminary Analysis (h-m1)}
\label{tab:sep_results}
\begin{tabular}{llll}
\toprule
\textbf{Condition} & \textbf{Mean SEP} & \textbf{N samples} & \textbf{vs.\ DPO} \\
\midrule
GRPO-binary & 0.2371 & 192 & $-$0.0006 (lower) \\
GRPO-error-type & 0.2371 & 192 & $-$0.0006 (lower) \\
DPO & 0.2377 & 189 & --- \\
\bottomrule
\end{tabular}
\end{table}

Mann-Whitney U (GRPO-binary vs.\ DPO): $U = 18{,}346.5$, $p = 0.4248$ (not significant; required $p < 0.05$).
Effect size: $-0.0072$ (GRPO slightly \emph{lower} than DPO, opposite of hypothesis direction).

\textit{Note:} GRPO-binary and GRPO-error-type produce identical SEP statistics because the h-m1 analysis
reused h-e1 checkpoints, and the checkpoint aliasing confound (Section~\ref{sec:aliasing}) caused
both conditions to analyze the same aliased checkpoint-100 files.

Figure~\ref{fig:gate_sep} shows the SEP distributions for GRPO and DPO side by side.
The distributions are nearly overlapping, with medians of approximately 0.237 for
both conditions. There is no visual or statistical evidence of SEP superiority for GRPO.

\textbf{The raw/proportion dissociation.} This result stands in striking contrast
to the proof-of-concept raw edit distance finding: GRPO produces
+250\% higher \emph{absolute} semantic AST edit distance (h-e1 synthetic PoC data),
yet the \emph{proportion} of edits targeting semantic nodes is essentially identical
($\approx$0.237 for both; h-m1 preliminary, $n_\text{eff}\approx2$, underpowered).
Figure~\ref{fig:ast_edit_dist} illustrates the distribution of edit distances by node category,
revealing that both methods increase semantic and surface edits proportionally.

A confirmatory dry-run smoke test on 10 fresh HumanEval+ problems (3 matched
pairs) also showed SEP $\approx$ 0.238 for both GRPO and DPO, providing weak
corroborating evidence that the near-equality is not an artifact of the
main analysis, even before considering the checkpoint aliasing confound.

\subsection{Checkpoint Aliasing Confound (RQ3)}
\label{sec:aliasing}

\textbf{Critical finding:} The h-m1 analysis was severely compromised by checkpoint
aliasing, a confound not, to our knowledge, previously documented in the RL fine-tuning literature.

\textbf{What happened:} The h-e1 proof-of-concept was a smoke test that saved only
2 real GRPO checkpoints (step-100 and step-200). The h-m1 analysis was designed
to use checkpoints from steps 100--1000 for its 27-pair KL-matched analysis.
However, steps 300--1000 fell back to checkpoint-100 (the earliest saved
checkpoint) when the requested checkpoint files were not found.

\textbf{Consequence:} 25 of 27 analysis pairs aliased to checkpoint-100. The
192 GRPO SEP values were derived from effectively 2 distinct checkpoints
(not 27 independent measurements), violating the independence assumption of
the Mann-Whitney test. The Spearman correlation was undefined (NaN) due to
zero variance in the x-axis across aliased pairs. The nominal $n=192$ collapses
to $n_\text{eff} \approx 2$.

Figure~\ref{fig:sep_kl_traj} illustrates the aliasing directly: the
SEP-vs-KL trajectory shows a near-flat line with almost no variation across
the 27 nominal pairs. Figure~\ref{fig:reward_scatter} shows the tight clustering of reward values
across checkpoint pairs, further evidence of aliasing.

\textbf{The confound is invisible from pass@1 metrics.} A team relying only on
final benchmark scores would not detect this aliasing. The structural efficiency
framework, by requiring checkpoint diversity verification, surfaces this failure
mode explicitly.

\textbf{Recommended safeguard:}
\begin{verbatim}
unique_checkpoints = len(set(checkpoint_paths))
assert unique_checkpoints >= N_min, (
    f"Checkpoint aliasing detected: only {unique_checkpoints} unique "
    f"checkpoints found (required >= {N_min}). Abort analysis."
)
\end{verbatim}

We recommend $N_\text{min} = 10$ for mechanism-level analysis and $N_\text{min} = 5$
for existence-level checks.
